\section{与文献的关系}\label{sec:literature}

\paragraph{戒烟、理性成瘾与目标相对恢复。}
戒烟构成全文贯穿的方案环境。\citet{BeckerMurphy1988} 用成瘾存量和相邻互补性刻画稳定的前瞻偏好，\citet{BeckerGrossmanMurphy1994} 则把这一结构用于香烟需求。后续研究讨论变化的短视程度 \citep{OrphanidesZervos1998}、现在偏误与自我控制需求 \citep{GruberKoszegi2001}、戒断与戒烟决策 \citep{SuranovicGoldfarbLeonard1999}、线索触发选择 \citep{BernheimRangel2004}，以及强迫性消费中的反复 \citep{Winston1980}。本文的通行次序隔离了一个互补对象：在给定上述模型所表示的其他后果条件下，个体如何对有向调整记录本身进行排序。

预先给定但不具强制约束力的目标，可以理解为治疗参照点或分阶段减量平台。目标可以支持自我调节 \citep{Hsiaw2013,KochNafziger2011}，承诺契约也已经直接用于戒烟研究 \citep{GineKarlanZinman2010}。个体可以未达到目标、越过目标或远离目标。内生目标形成、愿望水平修订和承诺机制设计属于应用层面的进一步扩展。

\paragraph{参照依赖与目标型状态效用。}
参照依赖模型相对于一个参照点评价结果，并允许参照点两侧具有不对称的敏感度 \citep{TverskyKahneman1991,KoszegiRabin2006}。本文表示中最接近这类模型的是仅依赖端点的子类。在任意固定的正则相容支撑赋值下，当 $\kappa_L=\kappa_R=0$ 且 $\alpha_L=\alpha_R=0$ 时，通行价值退化为
\[
 U_\Lambda(\gamma)=P_\Lambda(\gamma(1))-P_\Lambda(\gamma(0)).
\]
在初始状态固定时，每个支撑赋值于是只依赖以目标为中心的终止状态。一般表示在支撑赋值层面额外加入两种路径价值来源：由有向跨越次数产生的总通行调整，以及由实际目标达成产生的边界贡献。本文维持目标预先给定的环境，把基于预期形成参照点的问题留给具体应用。

\paragraph{序列、经历与程序价值。}
关于序列和延展经历的偏好可见于 \citet{LoewensteinPrelec1993} 与 \citet{VareyKahneman1992}。程序效用允许结果的生成方式本身进入福利 \citep{FreyBenzStutzer2004,FreyStutzer2005}；\citet{MinardiSavochkin2024} 则公理化了过去消费的可记忆效应。本文的通行分量保留一侧、方向、状态水平与重复次数，同时忽略日历时间、持续时间、新近性、适应以及彼此独立的闭合经历之间的顺序。

形如 $g(c,A)=\int_c^A q(s)\,\dd s$ 的参数积分为一次单调调整所跨越的状态水平赋值。有向通行测度把这一思路以序数、非参数方式扩展到重复的非单调历史、目标两侧、两个方向以及目标达成原子。

\paragraph{可加测量与受限定义域。}
从有限消去条件到可加表示的路线是经典的；参见 \citet{Scott1964}、\citet{KrantzEtAl1971} 和 \citet{Wakker1989}。\citet{KraftPrattSeidenberg1959} 在有限定性概率中发现的障碍，与本文使用平衡的有限证书而不只依赖成对一致性的做法非常接近。\citet{Fishburn2001} 研究有限乘积域上所需的不同比较对数量，同时允许证书中出现整数重复次数。本文的算术构造固定六个账本列，并让一个路径可实现族中的原始系数重复次数发散。

部分定义的运算是广延测量中的核心问题。\citet{Luce1967} 在不交并仅部分定义时得到有限可加概率表示，\citet{Mackenzie2021} 则从原子聚集条件推导相关假设。本文通过一个通用通行群对可行路径拼接作代数完备化，但只对实际比较所生成的揭示锥进行排序。受限域上的可加表示见 \citet{ChateauneufWakker1993}、\citet{Kopylov2007} 和 \citet{QinRommeswinkel2022}。允许形式差分的次序保持偏序，使实数定价步骤更接近多效用表示，而不是全局 H\"older 标量化；参见 \citet{Ok2002}、\citet{DubraMaccheroniOk2004} 和 \citet{EvrenOk2011}。

\paragraph{连续性、原子、路径与原始平衡。}
在定性概率中，单调连续性是从有限可加性走向可数可加性的经典路线 \citep{Villegas1964,ChateauneufJaffray1984,Chateauneuf1985}；相关连续性限制也被用于正则化先验族 \citep{ChateauneufMaccheroniMarinacciTallon2005}。本文的连续性条件由路径几何生成：只有几何交集为空的非目标穿孔区间被视为可忽略，而到达目标的序列可以保留一个单点残差。

忽略顺序而保留重复次数的汇总与 Parikh 向量有关 \citep{Parikh1966}，而路径签名则保留非交换的顺序信息 \citep{Chen1954,HamblyLyons2010}。纤维定理把 Euler 迹变换 \citep{Abrham1980,FleischnerEtAl1992} 专门化到带标签的一维路径，并固定端点不平衡与目标达成标签。总水平跨越计数有经典分析传统 \citep{Lochowski2017}；在本文中，它们被用于偏好表示以及端点--总通行分解。

回路、共形不可分解性和 Graver 基都是标准概念 \citep{Graver1975}。相关复杂度结果包括 \citet{BersteinOnn2009} 与 \citet{TatakisThoma2015}。本文的原始元素本身就是一个回路，对每个参数值都只有六列，其增长来自可达路径账户中的算术系数。总通行调整是一个静态、对反转保持偶对称的分量。进入--退出不可逆性 \citep{Dixit1989}、动态迟滞 \citep{MielkeRoubicek2015}，以及博弈和排序中的势--循环分解 \citep{CandoganEtAl2011,JiangEtAl2011}，提供了相邻但不同的结构。
